% hackweek_snowpack preamble — a projected 16:9 talk. The theme is picked in
% colors.tex; nothing here names a palette colour.
%
% Architecture follows ~/DavidFillmore/Media/decks/agentic_ai (16:9 beamer,
% Tectonic/XeTeX, fontspec before colors.tex, listings not minted), NOT that
% repo's templates/nasa_blue beamerposter build: there is no poster wrapper here, so the 11x8.5in page,
% tiles.tex/poster.tex and the banner-headroom machinery are all absent along
% with the type-scale trouble they cause.
%
% Berkeley is deliberately NOT used either, though agentic_ai uses it. This
% deck is half figures, and Berkeley's section sidebar costs ~1.6cm of every
% slide's width — which is exactly the width a full-bleed plot wants. The
% chrome here is one NASA-blue frametitle bar and a frame number: the
% "minimal feel" the deck was asked for.
\documentclass[aspectratio=169]{beamer}
\usetheme{default}
\useinnertheme{default}
\useoutertheme{default}

% fontspec is loaded ONLY so the CLI transcript panel can use Menlo, and it
% must come BEFORE colors.tex (which loads helvet). Two traps, both silent:
% fontspec resets \rmdefault/\sfdefault to Latin Modern, so loaded after
% colors.tex it takes the whole deck off Helvetica; and it switches
% \encodingdefault from T1 to TU, while the deck's Helvetica exists only as
% T1/phv — every request then falls back to LM Roman and the deck turns serif
% with no error and a clean build. Restore T1 here; \clifont carries its own
% TU encoding from \newfontfamily.
\usepackage{fontspec}
\renewcommand{\encodingdefault}{T1}\normalfont

\usepackage[T1]{fontenc}
\usepackage{lmodern}

% Tectonic reads UTF-8 natively, so a literal em dash has no mapping into the
% T1-encoded body font and is silently dropped ("Missing character"). inputenc
% is a no-op on a unicode-native engine, so map it by catcode.
\catcode`\—=\active
\def—{\textemdash}

\usepackage{graphicx}
\usepackage{tikz}
\usetikzlibrary{shapes.geometric, arrows.meta, positioning, fit, backgrounds,
                calc, decorations.pathreplacing, shapes.callouts, shapes.symbols}
\usepackage{booktabs}
\usepackage{amsmath}
\usepackage{amssymb}

% listings, never minted — no shell-escape under Tectonic.
\usepackage{listings}
\definecolor{lstbg}{RGB}{245,245,245}
\lstset{basicstyle=\footnotesize\ttfamily, breaklines=true, frame=single,
        backgroundcolor=\color{lstbg}}

% Palette, symbolic seam, logo binding and graphics paths.
% The theme switch, and everything shared across themes.
%
% ---- SWITCH THE DECK'S THEME BY EDITING THE ONE LINE BELOW ----------------
% themes/ holds one file per theme: the palette, the symbolic bindings the
% slides are written against, and \themelogo. Nothing else in the deck names a
% palette colour, so this line is the whole change.
% anthropic_orange — palette and symbolic bindings. \input by colors.tex.
%
% After ~/DavidFillmore/Media/templates/anthropic_orange/colors.tex, with ONE
% deliberate departure. That template is single-seam: it sets
% accentSoft = accent, so every tier is the one clay. This deck cannot use
% that. Slides 02 and 04 are built on a real contrast between a block and an
% alertblock — specification against constraint, archive against trap — and
% under a single seam the two blocks become the same colour and the argument
% goes with them. So accentSoft is bound to Anthropic's dark brown, the other
% end of the same brand family: the alert tier reads as the grave one rather
% than the loud one, which suits a warm palette with no red in it.

% --- Anthropic palette ---
% Clay is extracted from Claude_by_Anthropic.svg; the browns are the family it
% sits in, as recorded in Media/decks/agentic_ai/colors.tex.
\definecolor{claudeClay}{RGB}{217,119,87}    % Anthropic orange
\definecolor{claudeBrown}{HTML}{3D3929}      % dark brown
\definecolor{claudeKraft}{HTML}{D4A27F}      % kraft, the light end
\definecolor{claudeBrownMid}{HTML}{886D54}   % midpoint of brown and kraft

% --- Symbolic seam ---
% Slides use ONLY these; never the claude* names above.
%
% CONTRAST, measured, not assumed: white on clay is 3.12:1. That clears WCAG AA
% for large text and does not clear it for body text, so clay is a GROUND for
% bold headings only — the frametitle bar and block titles — and never carries
% body copy. Body-weight text takes accentDeep (11.9:1 on white) or accentSoft
% (11.0:1). This is the same trade the shipped anthropic_orange template and
% agentic_ai's \claudeslide both make, and it is fine on a projector at size;
% it is written down here so it stays a decision rather than an oversight.
\colorlet{accent}{claudeClay}                 % primary chrome
\colorlet{accentSoft}{claudeBrown}            % second tier: alerts, annotation
\colorlet{accentDeep}{claudeClay!62!black}    % deep ink: diagram text, captions
\colorlet{accentFill}{claudeKraft!32!white}   % diagram node fill
\colorlet{accentTint}{claudeClay!10!white}    % block body ground
\colorlet{mutedInk}{claudeBrown!52!white}     % de-emphasised labels

\newcommand{\themelogo}{Claude_symbol_clay.pdf}

% % nasa_blue — palette and symbolic bindings. \input by colors.tex.
% After ~/DavidFillmore/Media/templates/nasa_blue/colors.tex, extended with the
% three seam names this deck's slides need beyond the shipped accent trio.

% --- Official NASA brand palette ---
% NASA's core brand is three colours (blue, red, white), stable since 1992.
% NASA Blue PMS 287 C, NASA Red PMS Bright Red C. White is the third — use
% plain white.
\definecolor{nasaBlue}{HTML}{0B3D91}
\definecolor{nasaRed}{HTML}{FC3D21}

% --- Symbolic seam ---
% Slides use ONLY these; never the nasa* names above. That is the portability
% contract that makes a theme switch a one-line edit in colors.tex.
\colorlet{accent}{nasaBlue}              % primary chrome: title bar, block titles
\colorlet{accentSoft}{nasaRed}           % second tier: alerts, annotation, gates
\colorlet{accentDeep}{nasaBlue!70!black} % deep ink: diagram text, captions
\colorlet{accentFill}{nasaBlue!16!white} % diagram node fill
\colorlet{accentTint}{nasaBlue!8!white}  % block body ground
\colorlet{mutedInk}{black!42}            % de-emphasised labels

% The meatball is a full-colour insignia and does not reverse to white by a
% fill swap; on a dark ground it needs a white plate.
\newcommand{\themelogo}{nasa_logo.pdf}

% --------------------------------------------------------------------------
%
% The seam the slides may use, and the only one:
%   accent      primary chrome — frametitle bar, block titles. Bold text on it
%               only; see the contrast note in themes/anthropic_orange.tex.
%   accentSoft  second tier — alertblocks, annotation, review gates
%   accentDeep  deep ink — diagram labels, \figcaption, body-weight emphasis
%   accentFill  diagram node fill
%   accentTint  block body ground
%   mutedInk    de-emphasised labels (artifact counts under the phase chain)
%   \themelogo  the theme's own mark
%
% A slide that reaches past these for a literal colour name has broken the
% contract and will not survive the next switch. Check with:
%   grep -ohE 'nasa[A-Za-z]+|claude[A-Za-z]+' frames/*.tex

% --- Vendor inks, deliberately NOT part of the seam -----------------------
% The CLI transcript panel paints Claude Code's own interface: the clay `>`
% prompt and the clay bullet are that product's colours, the way the NASA
% meatball is NASA's. They are content, not chrome, and must not follow the
% theme — under nasa_blue the panel still has to look like the terminal it is.
% Defined here rather than in a theme file for exactly that reason.
% (anthropic_orange defines claudeClay too, to the same value; a second
% \definecolor of a name already defined is a redefinition, not an error.)
\definecolor{claudeClay}{RGB}{217,119,87}

% --- Font binding ---------------------------------------------------------
% Helvetica (helvet is the shipped stand-in), kept across both themes even
% though Media's anthropic_orange template specifies Latin Modern Sans. Two
% reasons: LM Sans sets about 15% narrower, so every hand-tuned column width,
% figure height and \vspace in frames/ would need re-measuring; and the
% plots carry their own Helvetica-like sans, which a Latin Modern deck would
% sit oddly against. \familydefault is set explicitly because nothing else
% picks sans by default.
\usepackage{helvet}
\renewcommand{\familydefault}{\sfdefault}

% --- Graphics paths (relative to slides/) ---------------------------------
% Both asset directories live INSIDE the deck. The marks were copied out of
% ~/DavidFillmore/Media/graphics/logos/ rather than referenced across the
% filesystem, so this repo builds the deck on its own: a relative path out to
% another checkout would break for anyone who cloned only earth-eval.
% Provenance and usage rules: graphics/logos/SOURCES.md.
\graphicspath{%
  {figures/}%
  {graphics/logos/}%
}


% --- Beamer chrome ------------------------------------------------------
% Written against the symbolic seam (accent / accentSoft / accentDeep) rather
% than the literal nasa* names, so the deck stays portable to another theme's
% colors.tex — the same contract templates/ documents.
\setbeamercolor{structure}{fg=accent}
\setbeamercolor{palette primary}{bg=accent, fg=white}
\setbeamercolor{frametitle}{bg=accent, fg=white}
\setbeamercolor{title}{fg=accent}
\setbeamercolor{block title}{bg=accent, fg=white}
\setbeamercolor{block body}{bg=accentTint, fg=black}
\setbeamercolor{block title alerted}{bg=accentSoft, fg=white}
\setbeamercolor{block body alerted}{bg=accentSoft!8!white, fg=black}
\setbeamercolor{item}{fg=accent}
\setbeamercolor{subitem}{fg=accentSoft}
\setbeamerfont{block title}{series=\bfseries, size=\small}
\setbeamertemplate{navigation symbols}{}
\setbeamertemplate{blocks}[rounded][shadow=false]
\setbeamertemplate{itemize item}{\textbullet}
\setbeamertemplate{itemize subitem}{--}

% Wider text block than the default outer theme's 1cm: this deck's figures are
% the argument and every millimetre of width shows up in them.
\setbeamersize{text margin left=0.55cm, text margin right=0.55cm}

% Full-bleed frametitle bar. The stock `default' outer theme sets the title as
% plain text with no ground, so the bar has to be built here; a beamercolorbox
% at \paperwidth is what makes it run edge to edge rather than stopping at the
% text margins.
\setbeamerfont{frametitle}{series=\bfseries, size=\large}
\setbeamertemplate{frametitle}{%
  \nointerlineskip%
  \begin{beamercolorbox}[wd=\paperwidth, sep=0pt,
                         leftskip=0.55cm, rightskip=0.55cm]{frametitle}%
    \vskip0.42em\usebeamerfont{frametitle}\insertframetitle\vskip0.46em%
  \end{beamercolorbox}%
}

% Footline: EMPTY, at David's direction. It carried the frame number alone in
% the accent at the right margin; no bar and no author/title repeat, because a
% nine-slide deck earns no persistent strip — and on inspection the number did
% not earn its corner either.
%
% An empty template is zero-height, so every frame gains the few millimetres the
% number occupied. That can only relieve an Overfull \vbox, never cause one, but
% it does mean a frame tuned against the old body height now sits slightly high.
% Restore by putting \insertframenumber back in the group below.
\setbeamertemplate{footline}{}

% --- \fitline{<content>} -------------------------------------------------
% Set <content> on one line within \linewidth, shrinking only if it would
% otherwise wrap. Never enlarges.
\newlength{\fitlinewd}
\newcommand{\fitline}[1]{%
  \settowidth{\fitlinewd}{#1}%
  \ifdim\fitlinewd>\linewidth
    \resizebox{\linewidth}{!}{#1}%
  \else
    #1%
  \fi
}

% --- Link convention (deck-wide) -----------------------------------------
% EVERY link renders as REGULAR TEXT, UNDERLINED — no color shift, no
% italics, no beamer link box. Routed through a macro so the convention is
% enforced by construction: change it here, not at each call site. \mbox keeps
% a URL from breaking mid-rule at a line end.
\newcommand{\decklink}[2]{\href{#1}{\mbox{\underline{#2}}}}
\hypersetup{colorlinks=false, pdfborder={0 0 0}}

% \codepath{...} — typewriter formatting for path and identifier references.
\newcommand{\codepath}[1]{\texttt{\small #1}}

% --- Caption lines -------------------------------------------------------
% \figcaption{...} — the one line of prose a figure slide is allowed. Set
% small, in the deep tier, so it reads as annotation rather than body text.
\newcommand{\figcaption}[1]{%
  {\centering\small\color{accentDeep}#1\par}}

% \figaside{...} — the qualifying line UNDER a caption: the caveat, the
% product name a number is meaningless without, the second reading.
%
% ITALIC, not merely a second colour. Under nasa_blue the tier was carried by
% hue alone — NASA red against NASA blue — and that broke on the switch to
% anthropic_orange, where accentSoft and accentDeep are two warm browns a few
% steps apart and the aside merged into the caption above it. A warm
% monochrome palette has no second hue to spend, so the distinction is made
% structurally and works under either theme.
% The optional argument is the alignment: \centering by default, because the
% usual place for an aside is under a centred figure caption. Slide 05 passes
% \raggedright — its aside closes a left-aligned text column, where a centred
% line reads as a mistake rather than as a distinct note.
\newcommand{\figaside}[2][\centering]{%
  {#1\scriptsize\itshape\color{accentSoft}#2\par}}

% --- Claude Code CLI transcript panel ------------------------------------
% Authentic terminal read: set in Menlo, the macOS Terminal font, loaded
% through fontspec (Tectonic runs XeTeX, so the system font is available by
% name). Content is real Unicode in a monospace grid, NOT TikZ-drawn glyphs —
% keeping every mark on the same character cell is what makes it read as a
% terminal. TikZ draws the panel rectangle only. Menlo lacks the CLI's own
% U+23FA and U+23BF, so U+25CF and U+2514 stand in.
\newfontfamily\clifont{Menlo}[Scale=0.90]

\definecolor{cliBg}{HTML}{1F1E1D}      % warm near-black terminal ground
\definecolor{cliBorder}{HTML}{3A3733}
\definecolor{cliFg}{HTML}{E8E6E3}
\definecolor{cliDim}{HTML}{8B8580}     % tool arguments, result lines

\newcommand{\cliprompt}[1]{{\color{claudeClay}\bfseries>}\ {\color{cliFg}#1}}
\newcommand{\cliassist}[1]{{\color{claudeClay}●}\ {\color{cliFg}#1}}
\newcommand{\clitool}[2]{{\color{claudeClay}●}\ {\color{cliFg}#1}{\color{cliDim}(#2)}}
\newcommand{\cliresult}[1]{\hspace{0.9em}{\color{cliDim}└\hspace{0.4em}#1}}

% \clipanel[<width>]{<lines>} — the transcript panel. Separate lines with \\.
\newcommand{\clipanel}[2][0.90\textwidth]{%
  \begin{center}
  \begin{tikzpicture}
    \node[draw=cliBorder, fill=cliBg, rounded corners=5pt, inner sep=11pt,
          text width=#1, align=left, font=\clifont\small, text=cliFg] {#2};
  \end{tikzpicture}
  \end{center}}
